\documentclass[11pt]{article}
\usepackage[T1]{fontenc}
\usepackage[utf8]{inputenc}
\usepackage{lmodern}
\usepackage[a4paper,margin=28mm]{geometry}
\usepackage{amsmath,amssymb,amsthm}
\usepackage{microtype}

\usepackage[hidelinks]{hyperref}
\usepackage{url}
\hypersetup{pdftitle={A compositional inverse for the support polynomials of comb preorder polytopes},pdfauthor={Xiaojian Zeng}}
\newtheorem{theorem}{Theorem}
\newtheorem{lemma}{Lemma}
\newcommand{\Q}{\mathbb{Q}}
\newcommand{\Z}{\mathbb{Z}}
\newcommand{\Cset}{\mathcal{C}}
\DeclareMathOperator{\supp}{supp}
\title{A compositional inverse for the support polynomials\protect\\ of comb preorder polytopes}
\author{Xiaojian Zeng\thanks{PBC School of Finance, Tsinghua University, 100190, China.\newline
Email: \href{mailto:xiaojianzeng@foxmail.com}{\texttt{xiaojianzeng@foxmail.com}}.}}
\date{September 14, 2026}

\begin{document}
\maketitle

\begin{abstract}
We prove the support-refined compositional inverse proposed by Athanasiadis
and Chapoton for comb preorder polytopes. All lower-ideal inequalities reduce
to prefix inequalities for nonnegative integer costs, with local support
weights $w_0=1+t$ and $w_c=2t+ct^2$ for $c\geq1$. A postorder correspondence
identifies the resulting weighted words with nonempty ordered forests.
Accounting for their canonical first leaf gives the exact normalization and
yields the inverse $y/\bigl((1+y)(1+(1+t)y)(1+(1+t)y/t)\bigr)$.
The argument holds in $\Q(t)[[x]]$ with the support parameter indeterminate.
\end{abstract}

\section{The comb inverse}

Let $\tau_0$ be the empty poset. For $n\geq0$, form $\tau_{n+1}$ by taking
the disjoint union of $\tau_n$ with one singleton and then adjoining a new
minimum. Thus $\tau_n$ has $2n$ elements. Its preorder polytope is
\begin{equation}\label{eq:polytope}
 Q_n=\left\{a\in\mathbb{R}_{\geq0}^{\tau_n}:
       \sum_{v\in I}a_v\leq |I|
       \text{ for every lower ideal }I\text{ of }\tau_n\right\}.
\end{equation}
A lower ideal contains every predecessor of each of its elements.
We count lattice points by their number of positive coordinates:
\begin{equation}\label{eq:h}
 h_n(t)=\sum_{a\in Q_n\cap\Z^{\tau_n}}
                t^{|\supp(a)|},\qquad h_0(t)=1.
\end{equation}
These sums are finite: the ideal $\tau_n$ bounds the sum of all coordinates
by $2n$. The polynomial in \eqref{eq:h} is the raw support enumerator,
called the $h$-polynomial in Section~5 of \cite{AC}; it is distinct from the
Ehrhart $h^*$-polynomial of $Q_n$.

Athanasiadis and Chapoton's final unnumbered proposal in
Section~9.6 of \cite{AC} concerns the series
\begin{align}
 F(t,x)&=\sum_{n\geq0}t^{-n}h_n(t)x^{n+1},\label{eq:F}\\
 R(t,x)&=\sum_{m\geq1}(-1)^{m-1}
 \left(1+(1+t)^m t^{1-m}\sum_{j=0}^{m-2}t^j\right)x^m.
 \label{eq:R}
\end{align}
The inner sum is zero when $m=1$. Throughout, $t$ is an indeterminate.

\begin{theorem}\label{thm:main}
For the comb posets and support polynomials above,
\begin{equation}\label{eq:inverse}
 R(t,F(t,x))=x=F(t,R(t,x))\qquad\text{in }\Q(t)[[x]].
\end{equation}
In particular, the series \eqref{eq:R} is the rational series
\begin{equation}\label{eq:rational}
 R(t,x)=\frac{x}{(1+x)(1+(1+t)x)(1+(1+t)x/t)}.
\end{equation}
\end{theorem}

Chapoton's arbor recursions already enumerate lattice points by support and
height (Propositions~2.8--2.9 of \cite{Chapoton}). General invariance of preorder
support polynomials under order duality follows from Wang, Hou, Thawinrak,
Liu and Dai \cite{Wang} (Theorem~1.3 and the following paragraph); that result
settles Conjecture~5.4 of \cite{AC}. The proposal addressed here is the
separate explicit comb inverse in Section~9.6. Our proof works directly
with the minimum-rooted comb and does not require the duality theorem.
Tree equations and their inversion are classical; see, for example,
Banderier and Flajolet \cite{BF}, Example~3, equation~(24).
The contribution here is the reduction of the comb's entire ideal system
to support-weighted prefix costs and the forest normalization that produces
the proposed inverse. No equivalent resolution was found in the primary
passages inspected through September~14, 2026; this is a bounded literature
assessment, without a claim of absolute priority.

\section{From lattice points to ordered forests}

Write $q=1+t$ and $K=\Q(t)$. Index the spine of $\tau_n$ from its minimum
as $r_1<\cdots<r_n$, and let $s_i$ be the leaf covering $r_i$.
The covering relations are exactly $r_i<r_{i+1}$ and $r_i<s_i$.
Indeed, each recursive step creates the new $r_1,s_1$ and increases all old
indices by one.

\begin{lemma}\label{lem:costs}
For a nonnegative integer vector with coordinates
$a_i=a_{r_i}$ and $b_i=a_{s_i}$, define
\[
 c_i=a_i+\max(b_i-1,0).
\]
Membership in $Q_n$ is equivalent to
\begin{equation}\label{eq:prefix}
 \sum_{i=1}^k c_i\leq k\qquad(1\leq k\leq n).
\end{equation}
If $\Cset_n$ denotes the set of nonnegative integer words satisfying
\eqref{eq:prefix}, then
\begin{equation}\label{eq:weighted}
 h_n(t)=\sum_{(c_1,\ldots,c_n)\in\Cset_n}\prod_{i=1}^n w_{c_i}(t),
 \qquad w_0=q,\quad w_c=2t+ct^2\quad(c\geq1).
\end{equation}
\end{lemma}

\begin{proof}
Every nonempty lower ideal has the unique form
\[
 I_{k,S}=\{r_1,\ldots,r_k\}\cup\{s_i:i\in S\},
 \qquad 1\leq k\leq n,\quad S\subseteq\{1,\ldots,k\}.
\]
The spine part must be an initial segment, and a leaf requires exactly its
parent and the preceding spine vertices. Conversely, every displayed set
is an ideal. Its inequality is
\[
 \sum_{i=1}^k a_i+\sum_{i\in S}(b_i-1)\leq k.
\]
Maximizing over $S$ gives \eqref{eq:prefix}, in both directions; the empty
ideal imposes no constraint.

For a fixed cost $c$, sum $t^{\mathrm{1}_{a>0}+\mathrm{1}_{b>0}}$ over pairs
$(a,b)$ with $a+\max(b-1,0)=c$. At cost zero the pairs are $(0,0),(0,1)$.
For $c\geq1$, the two pairs of support one are $(c,0),(0,c+1)$.
The support-two pairs are $(c,1)$ and $(c-b+1,b)$ for $2\leq b\leq c$,
giving exactly $c$ pairs. This proves the local weights. Once the costs
are fixed, the coordinate pairs can be chosen independently and their
supports add, proving \eqref{eq:weighted}. The empty word has weight one.
\end{proof}

The local degree series is therefore
\begin{equation}\label{eq:W}
 W(t,u)=\sum_{c\geq0}w_c(t)u^c
 =q+\frac{2tu}{1-u}+\frac{t^2u}{(1-u)^2}
 =\frac{(1+(t-1)u)(q-u)}{(1-u)^2}.
\end{equation}

An ordered rooted tree has an ordered list of children at each vertex;
an ordered forest is an ordered list of such trees. All vertices here are
unlabelled. In postorder, the children are visited in their order before
their parent, and components are visited in their order.

\begin{lemma}\label{lem:forest}
For every $n\geq0$, the words in $\Cset_n$ correspond bijectively to
nonempty ordered forests on $n+1$ vertices, with postorder outdegree word
\begin{equation}\label{eq:degrees}
 (0,c_1,\ldots,c_n).
\end{equation}
If a degree-$c$ vertex has weight $z w_c(t)$, and $T(z)$ enumerates ordered
rooted trees with these weights, then, for
$H(z)=\sum_{n\geq0}h_n(t)z^n$,
\begin{equation}\label{eq:forest}
 T=zW(t,T),\qquad qzH(z)=\frac{T(z)}{1-T(z)}.
\end{equation}
\end{lemma}

\begin{proof}
Process \eqref{eq:degrees} from left to right using a stack of completed
trees, kept as an ordered list. At an entry $d$, take the last $d$ trees
in their existing order as the children of a new root and replace them
by that tree. An entry zero simply appends a leaf; in particular the
initial zero creates one leaf.

Before processing $c_i$, the stack has
$i-\sum_{j<i}c_j$ trees: $i$ vertices and $\sum_{j<i}c_j$ edges have been
created. Hence the operation is legal exactly when
$\sum_{j\leq i}c_j\leq i$. After it the stack has
$i+1-\sum_{j\leq i}c_j\geq1$ trees. Inductively, concatenating the
postorder words of the current trees gives precisely the processed word.
The final stack is the required nonempty forest.

Conversely, a nonempty forest starts its postorder traversal at a leaf.
Processing any subtree by the stack rule leaves earlier stack entries
unchanged and appends exactly that subtree: this follows by induction on
its children. Thus the last $d$ completed trees at a degree-$d$ parent
are exactly its children, in order. Applying this observation component
by component reconstructs the original forest uniquely. Each operation
is legal, so its degree word satisfies \eqref{eq:prefix}. This establishes
both directions of the bijection.

Root decomposition now gives $T=z\sum_{c\geq0}w_cT^c$ and
$T=qz+O(z^2)$. Nonempty forests have series $T+T^2+\cdots=T/(1-T)$.
At each size there are only finitely many tree shapes, degrees and local
decorations, so these equations are valid coefficientwise even though the
degree set is unbounded. They define $T\in zK[[z]]$ uniquely by recursion.

The forest for a word in $\Cset_n$ has weight
\[
 z^{n+1}w_0\prod_{i=1}^n w_{c_i}
   =qz\,z^n\prod_{i=1}^n w_{c_i}.
\]
The extra vertex is the uniquely specified \emph{first} postorder leaf.
Its two cost-zero decorations contribute $1+t=q$; no leaf is chosen
arbitrarily, so no factor counting the leaves occurs. Summing proves the
second equation in \eqref{eq:forest}. For $n=0$, the single-vertex forest
has weight $qz$, giving the required constant term $H(0)=1$.
\end{proof}

\section{Formal elimination and the proposed coefficient series}

\begin{proof}[Proof of Theorem~\ref{thm:main}]
By definition $F(t,x)=xH(x/t)$. Set $U=T(x/t)$ and abbreviate $F=F(t,x)$.
Equation \eqref{eq:forest} gives
\begin{equation}\label{eq:U}
 F=\frac{tU}{q(1-U)},\qquad
 U=\frac{qF}{t+qF},\qquad U=\frac{x}{t}W(t,U).
\end{equation}
In the factorization \eqref{eq:W}, substitution yields
\[
 1-U=\frac{t}{t+qF},\qquad
 1+(t-1)U=\frac{t(1+qF)}{t+qF},\qquad
 q-U=\frac{qt(1+F)}{t+qF}.
\]
It follows that $W(t,U)=q(1+F)(1+qF)$, and therefore
\begin{equation}\label{eq:elimination}
 x=\frac{tU}{W(t,U)}
   =\frac{F}{(1+F)(1+qF)(1+qF/t)}.
\end{equation}

It remains to identify this rational expression with the coefficient
series actually proposed in \eqref{eq:R}. Using an auxiliary variable $y$,
the initial $1$ in each coefficient sums to $y/(1+y)$. For the remaining
terms put $m=k+2$, then $a=j$ and $b=k-j$. Their sum is
\begin{align*}
 -\frac{q^2y^2}{t}\sum_{k\geq0}(-qy)^k\sum_{j=0}^k t^{j-k}
 &=-\frac{q^2y^2}{t}\sum_{a,b\geq0}(-qy)^a(-qy/t)^b\\
 &=-\frac{q^2y^2}{(1+qy)(t+qy)}.
\end{align*}
Every fixed degree receives finitely many terms, justifying the
rearrangement as a formal double geometric sum. Consequently,
\begin{align*}
 R(t,y)&=\frac{y}{1+y}-\frac{q^2y^2}{(1+qy)(t+qy)}\\
 &=\frac{ty}{(1+y)(1+qy)(t+qy)}
 =\frac{y}{(1+y)(1+qy)(1+qy/t)},
\end{align*}
where the simplification uses
$(1+qy)(t+qy)-q^2y(1+y)=t$. This proves \eqref{eq:rational}, and
\eqref{eq:elimination} proves $R(t,F(t,x))=x$.

Both $F$ and $R$ have constant coefficient zero and linear coefficient
one. Such a series has a unique compositional inverse, obtained by solving
successive coefficients, so also $F(t,R(t,x))=x$.
All substitutions used zero-constant series; all denominators had nonzero
constant coefficient in $K$. In particular, $t$ and $q$ are units in this
field, and the division by $z$ implicit in \eqref{eq:forest} removes an
explicit factor. No analytic convergence or scalar specialization is used.
\end{proof}

The resulting support series also has a polynomial description. Substituting
$x=tz$ in \eqref{eq:elimination}, clearing denominators and cancelling $tz$
gives
\begin{equation}\label{eq:cubic}
 H=(1+tzH)(1+tqzH)(1+qzH).
\end{equation}
This identity holds in $\Z[t][[z]]$, and its constant-one solution is unique
by coefficient recursion. It includes the initial values
\[
 h_1(t)=1+3t+t^2,\qquad h_2(t)=1+8t+15t^2+8t^3+t^4.
\]
Unlike the intermediate quotient by $q$, \eqref{eq:cubic} can be specialized
at $t=-1$; it gives $H(z)=1/(1+z)$. At $t=0$ it gives $H(z)=1/(1-z)$.
These polynomial specializations do not define the Laurent series
\eqref{eq:F} at $t=0$.

\section{Exact checks and reproducibility}

The proof establishes every coefficient with $t$ indeterminate.
Supplementary programs use exact integers and Laurent-polynomial arithmetic
to check its combinatorial and algebraic steps. The verification design
uses different representations of the finite objects, as in the separate
integer-predicate computations of finite lattice configurations in
Zeng, Liu, Ratnavelu and Ong \cite{ZengChecks}, Sections~4 and~6.
That comparison concerns verification practice; their geometric results
are not used in the comb proof.

One program generates forests as recursive lists of trees. A separately
written program generates balanced-parenthesis encodings and reconstructs
the ideals directly from the minimum-adjunction recursion. The latter
checks the following finite ranges, with no numerical tolerances:
\begin{center}
\begin{tabular}{@{}ll@{}}

\textbf{Check} & \textbf{Coverage}\\

All-ideal versus prefix membership & $n=0,\ldots,5$; 198,627 vectors\\
Forests versus admissible cost words & $n=0,\ldots,8$; 6,917 forests\\
Canonical-leaf support normalization & $n=0,\ldots,8$\\
Local coordinate-pair weights & $c=0,\ldots,25$\\
Both formal compositions & through $x^{13}$ over $\Z[t,t^{-1}]$\\

\end{tabular}
\end{center}
The ambient lattice vectors in the first row are all nonnegative
$2n$-tuples of total at most $2n$; the forest comparison uses all 17,577
nonnegative length-$n$ words of total at most $n$ in its stated range.
The larger series check does not extend direct lattice enumeration beyond
$n=5$. These are finite consistency checks, not an extrapolation proof.

The supplementary file \texttt{comb\_checks.py} reproduces the latter
checks using only the Python~3 standard library; running
\texttt{python3 comb\_checks.py} writes \texttt{comb\_checks.json} alongside
the script. It is a publication adaptation of the separately written
checker, rather than an additional independent implementation.

\paragraph{Computational and AI assistance.}
OpenAI Codex assisted with proof development, literature searches, exact-check
programs and manuscript preparation. An additional model context examined
the argument and produced the separately written checker. Model output is
not a mathematical certificate: the all-degree justification is the
self-contained proof above. No proof assistant was used.


\small
\bibliographystyle{plain}
\bibliography{references}

\end{document}
